% !TeX program = xelatex

% 本文件由 md2tex.py 自动生成，请勿手改；请修改 Markdown 源后重新生成。

\documentclass[11pt,a4paper]{ctexart}

\usepackage{geometry}

\geometry{margin=2.5cm}

\usepackage{booktabs}

\usepackage{amsmath,amssymb}

\usepackage{graphicx}

\usepackage{hyperref}

\hypersetup{colorlinks=true,linkcolor=black,urlcolor=blue}

\setlength{\parindent}{2em}

\linespread{1.3}

\begin{document}

\title{主动式AGI的理论框架与落地约束：概念界定、持续认知闭环架构与开放研究问题}

\author{}

\date{\today}

\maketitle

\begin{abstract}


大语言模型智能体研究在过去三年中建立了从被动应答到任务自主执行的范式演进，但全部现有工作共享一个未被审视的前提：系统的认知过程由外部指令触发、随任务完成而终止。本文指出该前提构成通往通用人工智能的范式瓶颈，并提出第三种智能范式——主动式AGI（Persistent-Active AGI）：一种无需显式指令即可常驻运行、持续观测环境与自身状态、通过内部自省识别缺口、内生生成次级目标并在人机可控对齐下闭环迭代的通用智能形态。本文完成四项理论贡献：（1）建立三范式七维度的严格概念区分体系，以五条件判据界定主动式AGI，并论证当前全部主流Agent系统均不满足该判据，仅为"主动行为的模拟"；（2）提出五层持续认知闭环架构（FL-PCL），其中目标内生层引入经外源锚定改造的预期自由能框架，为无指令目标生成提供原则性候选判据；（3）系统论证内生目标漂移、自省不可靠、长期记忆腐化、常驻算力边界与无人值守对齐五大理论困境，揭示其共同根源为闭环的自指性、锚点稀缺性与能量非零性；（4）给出三级技术边界评估与极限资源推演，证明主动式AGI的本质边界在评价端而非执行端，顶配技术栈仅能实现"短切片全真、长跨度失真"的模拟主动性，当前商业可行形态止步于半常驻人机间歇耦合。本文最终提炼七条开放研究问题，构成从理论到受限实证的研究纲领。本文不声称实现主动式AGI，而是使其首次成为可检验、可分工、可接续的科学对象。

\par\noindent \textbf{关键词}：大语言模型；智能体；主动式AGI；持续认知闭环；自主目标生成；人机对齐
\end{abstract}

\section{1 绪论}

\subsection{1.1 研究背景}

自2022年末以来，以大规模预训练语言模型（Large Language Models, LLMs）为核心的智能系统经历了两轮范式演进。第一轮是\textbf{被动应答式AI}：以指令微调与人类反馈强化学习（RLHF）为技术主线，系统在用户显式提问后生成响应，无交互即处于休眠状态，其智能完全外生于人类输入。第二轮是\textbf{任务驱动型自主Agent}：以 ReAct（Yao et al., ICLR 2023）提出的"推理—行动"交替范式为起点，结合工具学习（Toolformer, Schick et al., NeurIPS 2023）、语言化自我反思（Reflexion, Shinn et al., NeurIPS 2023）与开放环境探索（Voyager, Wang et al., 2023）等机制，使系统能够在给定任务下进行多步自主规划与执行；AutoGPT、BabyAGI 等开源工程实践进一步将该范式推向公众视野。

然而，无论上述两轮范式在能力维度上取得多大进展，其\textbf{存在论前提}并未改变：系统的认知生命周期由外部指令触发、由任务完成而终止。换言之，当前主流Agent的"自主性"是\textbf{任务内的执行自主性}，而非\textbf{存在论意义上的持续自主性}。没有指令，就没有认知过程；没有任务，就没有目标结构。本文认为，这一局限构成了通往通用人工智能（AGI）道路上的一个根本性范式瓶颈：真实智能体的核心特征并非"接到任务后执行得好"，而是"在无任务时仍然在看、在想、在生成目标"。

据此，本文提出并系统性论证第三种范式——\textbf{主动式AGI（Persistent-Active AGI）}：一种具备常驻式持续认知闭环的通用智能范式，无需用户显式指令触发即可长期驻留运行，持续观测环境与自身状态，通过内部自省发现缺陷与缺口，内生地生成次级目标，并在人机可控对齐约束下完成规划、执行与记忆回写。

\subsection{1.2 国内外研究现状综述}

本节从五条与本文最相关的研究脉络梳理现状，并标注其与主动式AGI的关系。

\subsubsection{1.2.1 任务驱动型自主Agent}

ReAct（Yao et al., ICLR 2023）将思维链推理与外部环境行动交替交织，奠定了LLM Agent的基本回路结构；Toolformer（Schick et al., NeurIPS 2023）证明了模型可以自监督方式学习工具调用时机；Reflexion（Shinn et al., NeurIPS 2023）引入语言化的自我反思机制，使Agent能在失败后将口头反馈写入情节记忆以改进后续尝试；Voyager（Wang et al., 2023）在Minecraft开放环境中实现了由自动课程驱动的技能库持续积累，是现有工作中最接近"持续运行"形态的系统之一。Generative Agents（Park et al., UIST 2023）则展示了多智能体在沙盒社会中的记忆流、反思与规划机制，呈现出一定的"无显式任务下的行为生成"特征。

\textbf{与主动式AGI的关系}【现有技术可部分实现】多步规划、工具使用、反思与记忆机制均可复用为本文架构的执行层与自省层组件；但上述全部系统的顶层目标均由人类给定或由预设课程/脚本驱动，生命周期仍以任务或会话为单位，不具备无指令条件下的目标内生源。AutoGPT类系统在无有效外部约束时的目标漂移与空转问题，已在社区复现与评测中被广泛记录。

\subsubsection{1.2.2 自我进化智能体（Self-Evolving Agents）}

2025年起，自我进化方向出现系统性综述：Gao et al.（arXiv:2507.21046, TMLR 2026）以"进化什么、何时进化、如何进化"三维度组织该领域，覆盖模型、上下文、工具与架构四个进化对象；Fang et al.（arXiv:2508.07407）提出"系统输入—Agent系统—环境—优化器"的统一反馈回路框架，系统梳理了基于交互数据与环境反馈的自动增强技术。方法论层面，Self-Refine（Madaan et al., NeurIPS 2023）证明了"生成—反馈—精炼"的自循环可在单一模型内提升输出质量；基于自生成数据的训练式进化与基于经验的测试时进化均已成为活跃方向。

\textbf{与主动式AGI的关系}【现有技术可部分实现】+【当前学术界未解决开放难题】自我进化研究提供了"系统在无新指令下自我改进"的方法论雏形，是主动式AGI认知闭环的重要理论资源；但该脉络的进化信号（标量奖励、文本反馈、任务成功率）几乎全部来自外部给定的任务分布与评测器，"进化方向"本身并非内生。自我改进的稳定性、自省回路的可靠性（自修正退化、自我确认偏误）在现有文献中已被反复报告为未解决问题。

\subsubsection{1.2.3 主动推理与自由能原理（Active Inference / Free Energy Principle）}

理论神经科学一侧，Friston 提出的自由能原理与主动推理框架（Friston, 2010; Parr et al., 2022）将智能体建模为通过最小化预期自由能来持续更新世界模型并主动选择行动的自主系统，天然包含"无外部指令下的持续感知—行动循环"与内在动机（信息增益驱动的目标生成）。近年来已有将主动推理与深度学习、序列决策结合的计算化尝试。

\textbf{与主动式AGI的关系}【理论构想】主动推理在概念层面为"目标内生"与"常驻运行"提供了最成熟的规范性理论，是本文五层架构（尤其是目标内生层）的重要理论对照；但其与LLM级符号—语言认知系统的融合目前仅有零散探索，缺乏统一的、可验证的形式化桥接，属于理论构想而非可直接复用的技术。

\subsubsection{1.2.4 长时记忆与常驻记忆架构}

MemGPT（Paster et al., 2023）借鉴操作系统的虚拟内存分层思想，在有限上下文窗口之上构建主存—外存的分页管理机制，支持跨会话的长期状态维持；检索增强生成（RAG）与向量记忆库已成为Agent长期记忆的事实标准组件；Generative Agents 的记忆流与重要性评分、反思归纳机制提供了面向长期行为一致性的记忆组织方案。

\textbf{与主动式AGI的关系}【现有技术可部分实现】+【当前学术界未解决开放难题】记忆的读写、分层、检索机制可复用；但长期持续运行下的记忆累积失真、语义腐化、自我指涉污染（系统把自身的错误输出当作事实回写）与分布漂移尚无公认解决方案，现有研究多在数小时至数天的会话尺度上评估，与"常驻"所要求的月、年级尺度存在数量级差距。

\subsubsection{1.2.5 主动性研究：从主动对话到主动服务}

与本文最直接相关的是"主动性（proactivity）"研究脉络。Deng et al.（IJCAI 2023 Survey；扩展版 ACM TOIS 2025）首次系统综述了主动对话系统，将主动性形式化为 anticipation（预判）、initiative（发起）与 planning（规划）三要素，覆盖开放域、任务型与信息寻求型对话；其后续工作进一步提出 intelligence--adaptivity--civility 分类学，强调无节制的主动性会被感知为侵扰。Tang et al.（arXiv:2609.03727, 2026）将主动性从对话推进到"主动服务Agent"，提出以授权与风险约束的半马尔可夫决策过程来形式化"保持沉默、询问、协助、行动"的门控选择，并明确指出：长期记忆不是主动性的定义性条件，可靠的主动服务需要可校准的增量干预价值、可验证的授权与可恢复的执行。

\textbf{与主动式AGI的关系}【现有技术可部分实现】该脉络已在"何时不行动""干预时机校准""授权作为硬约束"等关键问题上给出了精细的决策论形式化，可直接为本文第五层（人机可控安全对齐层）提供理论支撑；但其研究对象仍是面向用户服务的干预决策，即"何时替用户做事"，而非"系统在无用户目标时的自我目标生成"，其沉默—干预的二元结构也不覆盖纯内生认知活动（自省、记忆整理、世界模型维护）。本文与该脉络的关系是继承与拓展：继承其授权—门控理论，拓展其对象域。

\subsection{1.3 研究缺口：本文的核心创新点}

综合上述五条脉络，可以精确刻画当前学术界的研究缺口：

\begin{enumerate}

  \item \textbf{自主性的层级缺口。} 现有研究覆盖了"任务内自主"（1.2.1）与"能力进化自主"（1.2.2），但对\textbf{存在论层面的持续自主}——无任务、无指令、无外部奖励信号条件下认知过程的自我维持与目标的自我生成——缺乏统一的定义、形式化模型与理论边界分析。主动推理（1.2.3）在规范性理论上最接近，但与LLM Agent技术栈脱节；主动性研究（1.2.5）在决策论上最精细，但研究对象是"面向用户的干预"，而非"面向系统自身的内生认知"。

  \item \textbf{目标来源的盲区。} 现有几乎全部Agent系统的目标来源可归结为三类：人类显式给定、预设课程/脚本隐式给定、外部奖励信号驱动。\textbf{由系统内部状态（知识缺口、记忆矛盾、能力缺陷、环境异常）触发、经自省回路内生生成次级目标的机制}，尚无可参照的架构级方案，也缺乏对目标漂移、价值偏移的理论刻画。

  \item \textbf{常驻运行的理论空白。} "持续运行"在现有文献中仅作为工程部署形态出现（如长驻服务进程），从未作为\textbf{理论对象}被研究：持续认知闭环应满足何种结构条件？其记忆、自省、目标生成三类子系统在长期运行下各自的失效模式是什么？算力与对齐约束给出何种可行性边界？本文首次将这些问题组织为一个完整的理论框架。

\end{enumerate}

\textbf{本文创新定位}：现有学术界研究的是"用户给定任务后的自主执行"；本文研究的是"\textbf{无任务、无指令下的持续自主认知与目标内生}"。本文不声称实现AGI，而是给出该范式的严格概念界定、五层持续认知闭环架构、五大核心理论困境的系统性论证，以及可供学界跟进的开放研究问题。

\subsection{1.4 本文主要研究工作与章节安排}

本文的贡献为纯理论性质，包含四部分：（i）构建"被动应答AI / 任务自主Agent / 主动式AGI"的三范式七维度严格区分体系（第2章）；（ii）提出主动式AGI的五层持续认知闭环理论架构，并逐层界定其定义、功能、输入输出、现有技术支撑与瓶颈（第3章）；（iii）系统论证无监督内生目标漂移、自省不可靠、长期记忆腐化、常驻推理算力边界、无人值守对齐安全五大理论困境（第4章）；（iv）给出落地可行性的三级分类评估、当代技术极限下的落地推演与七条开放研究问题（第5、6章）。第7章总结全文。

\section{2 核心概念辨析与范式对比}

\subsection{2.1 概念界定的必要性}

"主动（proactive）""自主（autonomous）""智能体（agent）"在当前文献中被高度混用：对话系统文献中的"主动性"指发起话轮的能力；Agent文献中的"自主性"指减少人工干预的执行能力；工程语境中的"常驻"仅指进程不退出。三者互不蕴含。若不先在存在论层面给出严格区分，任何关于"主动式AGI"的架构讨论都会退化为术语之争。本章给出一个三范式、七维度的区分体系，全部判据均为\textbf{可判定的结构性属性}（即给定系统描述即可客观归类），不诉诸意识、理解等不可测概念。

\subsection{2.2 三个核心概念的严格定义}

\textbf{定义1（被动应答式AI，Reactive-Responsive AI）。} 系统的全部计算活动构成对用户显式输入的函数响应：无输入则无任何计算状态转移（休眠）；有输入则产生输出后回到休眠。其形式特征为：（i）触发机制：唯一的外部指令触发；（ii）目标结构：无内部目标，仅有输入诱导的即时任务表征；（iii）生命周期：与单次请求—响应事务同构。典型实例：指令微调后的对话模型、无状态API式模型服务。

\textbf{定义2（任务驱动型自主Agent，Task-Driven Autonomous Agent）。} 系统满足：（i）触发机制：由人类给定的顶层任务（或预设课程/脚本隐式给定的任务序列）触发；（ii）在任务存续期内具备多步自主规划、工具调用、自我反思与中途修正能力，无需逐步人工干预；（iii）生命周期随任务达成、失败或资源耗尽而终止，任务终止后系统回到与定义1等价的休眠态。其自主性为\textbf{任务内的执行自主性}：自主的是"怎么做"，不是"做不做"，更不是"为何而做"。典型实例：ReAct/Toolformer 类单步回路系统、Reflexion 类反思系统、AutoGPT/BabyAGI 类工程框架、Voyager 类课程驱动探索系统。

\textbf{定义3（主动式AGI，Persistent-Active AGI，本文定义）。} 系统同时满足以下五个必要条件：

\begin{itemize}

  \item \textbf{(P1) 常驻运行}：系统的认知过程不依赖外部指令触发而持续存在，生命周期不以任务或会话为单位，休眠与唤醒由内部状态而非外部请求决定；

  \item \textbf{(P2) 持续观测}：系统在无任务状态下仍维持对环境与自身内部状态（记忆、能力、资源）的低频持续采样与表征更新；

  \item \textbf{(P3) 内部自省}：系统能对自身状态进行评价性分析，识别知识缺口、记忆矛盾、能力缺陷与优化空间；

  \item \textbf{(P4) 目标内生}：系统能仅由内部状态（而非外部指令或外部奖励信号）触发，生成、排序并维护一个次级目标队列；人类给定目标在此结构中作为高优先级外源目标进入同一队列，而非认知过程的唯一起点；

  \item \textbf{(P5) 可控对齐}：系统对行为按风险分级自治——低风险认知与行为自主迭代，高风险行为须经人类授权后方可执行，授权机制为硬约束而非软惩罚。

\end{itemize}

五条件联立为必要且（本文主张的）充分判据：缺P1退化为定义2；缺P4则任何"常驻"都退化为空转的守护进程；缺P5则系统不满足可部署性前提。

\subsection{2.3 七维度范式对比}

\textbf{表1 三类智能范式的结构性对比}

\begin{center}

\begin{tabular}{{p{1.6cm} p{3.4cm} p{4.1cm} p{4.4cm}}}

\toprule

对比维度 & 被动应答式AI & 任务驱动型自主Agent & 主动式AGI \\

\midrule

触发机制 & 用户显式输入，唯一触发源 & 人类给定顶层任务触发，任务内自主续行 & 内部状态触发为主，外部指令作为外源目标并入 \\

目标来源 & 无目标；输入诱导的即时任务表征 & 外生：人类显式给定或预设课程/脚本/奖励信号隐式给定 & 内生为主：自省回路识别缺口后生成次级目标；外生目标降格为队列中的高优先级项 \\

生命周期 & 与单次请求—响应事务同构 & 与任务同构：任务达成/失败/资源耗尽即终止 & 常驻：不以任务或会话为单位，跨任务连续存续 \\

运行状态 & 二值：响应中 / 休眠 & 二值：任务执行中 / 休眠 & 连续谱：观测、自省、目标维护、执行、记忆整理等多模态切换 \\

自省机制 & 无 & 有但任务内：反思仅服务于当前任务成败（如 Reflexion） & 有且任务外：自省指向系统自身的知识、记忆与能力结构 \\

终止条件 & 输出生成完毕 & 任务完成、失败或步数/预算耗尽 & 无自然终止；仅有外部停机、资源约束或对齐层的熔断终止 \\

智能本质 & 响应能力（response-ability） & 任务内的执行自主性（intra-task autonomy） & 存在论层面的持续自主性（persistent autonomy）：自我维持认知过程并自我生成目标 \\

\bottomrule

\end{tabular}

\end{center}

\subsection{2.4 关键论断：当前主流Agent均不属于主动式AGI}

依据表1的判据，可对第1章综述的代表性系统逐一归类论证：

\begin{enumerate}

  \item \textbf{AutoGPT/BabyAGI 类框架}：其"连续运行"是任务存续期内的循环执行，顶层目标由用户输入一次性给定，目标队列的源头是外生的；"连续"是循环语句意义上的连续，而非定义3中P1+P4意义上的常驻。任务队列耗尽后系统空转或停机，不满足P4。归类：定义2。

  \item \textbf{Voyager}：具备技能库跨任务积累与自动课程，是定义2的上界形态。但其自动课程由预设的总体目标（"发现尽可能多样的物品"）驱动，该目标为人类硬编码；探索的"内在动机"实为预设课程生成器的输出，非自省回路对内部缺口的识别。不满足P4。归类：定义2。

  \item \textbf{Generative Agents}：呈现出最丰富的无显式任务行为（自发社交、活动传播），看似最接近P4。但其行为生成器的顶层驱动是人工设定的人物种子身份与日程脚本，"自发行为"是种子设定的展开而非内生目标；且其运行依赖外部时钟逐 tick 驱动，无风险分级自治结构（P5缺失）。归类：定义2的高拟态形态——本文称之为"\textbf{主动行为的模拟}"（simulated proactivity），即以足够丰富的外生设定在外观上逼近内生主动性。这一区分是本章的重要结论：\textbf{外观上的主动行为不构成范式判据，目标来源的结构才是。}

  \item \textbf{主动服务Agent（Tang et al., 2026）一脉}：满足P2的部分要求（持续信号流上的状态估计）与P5的成熟形式化（授权门控、沉默的选项价值），但其决策对象是"何时介入用户事务"，服务机会的定义外生于用户；无用户目标时系统无认知活动，不满足P4。归类：定义2向定义3过渡带上、面向用户侧的中间形态。本文第五层将继承其授权—门控理论，并将其从"用户干预决策"拓展为"内生行为的风险分级"。

  \item \textbf{主动推理智能体}：在规范性理论上满足P1、P2与P4的雏形（预期自由能驱动下的持续感知—行动循环与信息增益目标），是唯一天然具备常驻结构的框架；但其现有计算实例均为小规模、特定域的生成模型系统，不具备LLM级的通用语言认知能力，且无风险分级对齐结构（P5缺失）。归类：理论层面的先驱形态，工程层面尚不构成定义3的实例。

\end{enumerate}

\textbf{小结}：截至本文写作时点，公开文献与公开系统中不存在满足P1--P5全部条件的系统；现有工作分别在单个条件上达到部分满足。这一事实既印证了第1章的研究缺口，也说明主动式AGI目前是\textbf{理论对象而非工程对象}——本文余下章节即按此定位展开：第3章给出使P1--P5结构化的五层架构，第4章论证其尚未被突破的理论障碍。

\section{3 主动式AGI五层持续认知闭环理论架构}

\subsection{3.0 架构总览与设计原则}

本章给出满足定义3中P1--P5条件的结构化模型：\textbf{五层持续认知闭环架构}（Five-Layer Persistent Cognitive Loop, 下称FL-PCL）。设计遵循三条原则：

\begin{enumerate}

  \item \textbf{闭环优先于流水线}。五层不是单向数据流，而是一个永续循环：观测\$\textbackslash{}to\$自省\$\textbackslash{}to\$目标\$\textbackslash{}to\$执行\$\textbackslash{}to\$（经对齐层过滤的）行动与记忆回写\$\textbackslash{}to\$新的内部状态\$\textbackslash{}to\$再观测。P1（常驻）由闭环的无终止迭代结构实现，而非由外部调度器实现。

  \item \textbf{目标的内生通道与外生通道并轨}。第三层维护统一目标队列，内生目标（源自第二层自省输出）与外生目标（源自人类指令）在同一优先级结构中被调度，从机制上落实P4中"外生目标降格为队列项"的判据。

  \item \textbf{对齐层为横切层而非末端层}。第五层不是执行完毕后的审查器，而是横切第三、四层的门控结构：目标入队需经风险分级，行动出栈需经授权校验，落实P5的硬约束属性。

\end{enumerate}

闭环的一次完整迭代记为一个\textbf{认知周期（cognitive cycle）}。与任务驱动Agent的"步（step）"不同，认知周期不以任务为单位度量，而是系统自我维持的基本时间单位；其频率由第二层估计的内部状态变化率自适应调节（状态平稳时降频以节约算力，状态剧变时升频），该设计同时为第4.4节的算力边界讨论埋下结构前提。

\subsection{3.1 第一层：环境持续观测层（Persistent Observation Layer）}

\textbf{定义。} 在无任务状态下，对环境信号流与系统内部状态进行低频、持续、带选择性注意采样的感知子系统。其输出不是"回答"，而是对\textbf{世界状态与自身状态的增量表征更新}。

\textbf{核心功能。} （i）外源观测：对授权范围内的外部环境信号（信息流、数据流、物理传感流，视部署域而定）进行事件化抽取；（ii）内源观测：对系统自身状态——记忆库存量与一致性指标、能力调用成功率的滑动统计、算力与资源预算——进行自检式采样；（iii）变化检测：将采样结果与既有状态表征比对，输出"显著变化事件"而非原始流，向上层传递的是差分而非全量。

\textbf{输入 / 输出。} 输入：授权环境信号流、内部状态探针读数、第二层下发的观测先验（即"当前认为值得注意什么"）。输出：增量状态表征 \$\textbackslash{}Delta\$s\_t 与显著变化事件队列。

\textbf{现有技术支撑。}【现有技术可部分实现】流式信息抽取、检索增强的状态跟踪、主动服务Agent一脉的状态与需求估计（Tang et al., 2026）、MemGPT式记忆分页提供的内源探针接口，均可复用为本层组件。内源观测中的"能力自检"可借用自我进化文献中的经验统计机制（Gao et al., 2025）。

\textbf{当前技术瓶颈。}【当前学术界未解决开放难题】（i）\textbf{无任务注意力的先验问题}：任务驱动系统的注意力由任务定义，无任务状态下"什么值得观测"缺乏原则性答案，本层只能依赖第二层下发的先验，而该先验的可靠性受制于自省质量（见4.2）；（ii）\textbf{观测成本—覆盖权衡}：持续采样的频率、广度与算力预算之间的最优分配无理论解；（iii）\textbf{内源观测的指涉可靠性}：系统对自身状态的表征本身就是模型生成的文本/向量，其作为"自检依据"的可信度未被任何形式化理论保证。

\subsection{3.2 第二层：内部认知自省层（Cognitive Introspection Layer）}

\textbf{定义。} 对第一层输出的状态表征进行评价性分析，识别系统自身的知识缺口、记忆矛盾、能力缺陷与优化空间的元认知子系统。这是FL-PCL区别于一切任务驱动架构的枢纽层：任务驱动Agent的反思以"任务是否完成"为评价基准，本层的自省以"\textbf{系统自身认知结构的完整性与一致性}"为评价基准。

\textbf{核心功能。} （i）缺口检测：在知识表征中定位高密度不确定区域（如检索反复失败的主题簇、内部问答自洽性低的断言集）；（ii）一致性审计：检测记忆库中的相互矛盾条目、过期条目与自我指涉污染（系统将自身先前错误输出当作事实回写形成的伪知识）；（iii）能力评估：基于第一层的能力统计数据，识别退化或欠发展的技能维度；（iv）缺口的价值评估：对检出的缺口按其重要性、可及性、风险进行初评，形成"自省报告"。

\textbf{输入 / 输出。} 输入：\$\textbackslash{}Delta\$s\_t、记忆库全文访问、能力统计、历史自省报告（防止重复发现）。输出：结构化自省报告——按价值初评排序的缺口清单，作为第三层目标生成的\textbf{唯一内生刺激源}。

\textbf{现有技术支撑。}【现有技术可部分实现】Reflexion 的语言化反思、Self-Refine 的自反馈循环、SelfCheckGPT 式的自洽性幻觉检测（Manakul et al., EMNLP 2023）、RAG系统中的冲突检索，分别提供了缺口检测与一致性审计的局部机制。Generative Agents 的反思归纳是"从记忆流提炼高层洞察"的先例。

\textbf{当前技术瓶颈。}【当前学术界未解决开放难题】（i）\textbf{自省的不可靠性}：LLM的自我评价能力被多项研究系统性质疑——自我修正在外部验证缺失时经常退化甚至负优化（自修正退化问题，见4.2的完整论证）；（ii）\textbf{评价基准的内源困境}：任务内反思有任务成败作为外部锚点，本层自省的评价基准是系统自己生成的，存在循环论证风险——用不可信的模型评价不可信的模型；（iii）\textbf{自省深度的终止问题}：对"我的自省是否可靠"的元自省可以无限递归，架构上必须设定强制截断深度，而截断点的选择缺乏原则。

\subsection{3.3 第三层：自主目标内生生成层（Endogenous Goal Generation Layer）}

\textbf{定义。} 将第二层输出的缺口清单转化为可执行、可排序、可撤销的次级目标，并与人类外生目标在统一队列中协同调度的子系统。本层是P4判据的机制实现，是本文架构的最高创新点。

\textbf{核心功能。} （i）目标合成：将缺口（陈述性描述）编译为具备前置条件、成功判据、资源预算与风险初级的目标对象；（ii）目标排序：维护统一优先级队列，排序键为内生价值评分与外生优先级的合成；（iii）目标生命周期管理：目标的挂起、合并、分解与撤销（当记忆审计表明目标基于伪知识时）；（iv）外生目标并入：人类指令作为最高优先级目标入队，与内生目标共享同一调度结构。

\textbf{候选形式化框架：预期自由能的改造引入。} 目标内生"生成什么目标"需要一个原则性判据，否则目标合成退化为启发式模板。本文引入主动推理的\textbf{预期自由能（Expected Free Energy, EFE）}作为候选规范性框架，并加以改造。经典EFE将策略 \$\textbackslash{}pi\$ 的价值写为：

\[G(\pi) = -\underbrace{\mathbb{E}_{Q(s|o,\pi)}[\ln P(o|s)]}_{\text{认识价值（信息增益）}} - \underbrace{\mathbb{E}_{Q(o|\pi)}[\ln P(o|C)]}_{\text{实用价值（偏好对齐）}}\]

即好策略同时最大化"减少世界模型不确定性"与"达成先验偏好"。迁移至本层：缺口 g 被编译为目标 \$\textbackslash{}pi\$\_g 的优先级正比于改造后的 Ĝ(\$\textbackslash{}pi\$\_g)，其中认识价值项对应"消除该缺口带来的内部表征不确定性下降"——这与自省层的缺口检测天然衔接；实用价值项中的偏好分布 P(o|C) 不再取自智能体自身先验，而\textbf{由人类价值观与系统章程外部锚定}（记为 C\_human），这是本文对经典EFE的关键改造：内生目标的"值得做"判据必须外源锚定，否则目标生成将滑向纯信息增益驱动的无约束探索（目标漂移的形式化起点，见4.1）。

【理论构想】需要明确：EFE在主动推理文献中针对的是低维连续状态空间的生成模型智能体，将其认识价值项定义在LLM的符号—分布式混合表征（知识库、向量记忆）之上，目前没有任何已发表的严格构造。本文给出的 Ĝ 是一个\textbf{框架性提案}而非已验证机制：其每一项的可计算近似（如以检索自洽性方差近似表征不确定性）都是独立的开放问题（见第6章）。

\textbf{输入 / 输出。} 输入：自省报告（内生刺激）、人类指令（外生刺激）、C\_human 锚定偏好、资源预算。输出：统一目标队列（含每个目标的成功判据、预算与风险初级）。

\textbf{现有技术支撑。}【现有技术可部分实现】目标分解与任务规划（ReAct后续工作、分层规划）、AutoGPT式目标队列工程结构、内在动机/好奇心驱动的探索（强化学习文献中的ICM、RND类机制提供了"信息增益作为目标信号"的实验先例）。

\textbf{当前技术瓶颈。}【当前学术界未解决开放难题】（i）\textbf{内生价值的可计算性}：Ĝ 各项在LLM表征上的可计算近似不存在公认方案；（ii）\textbf{目标膨胀与自我繁殖}：内生目标会生成子目标，无约束时目标队列呈指数膨胀，而"何时停止分解"无原则性判据；（iii）\textbf{内生—外生目标的冲突仲裁}：当内生目标与人类隐式偏好冲突（而非显式指令冲突）时的仲裁规则，涉及价值对齐的核心困难。

\subsection{3.4 第四层：任务规划与执行层（Planning \& Execution Layer）}

\textbf{定义。} 从统一目标队列取出目标，进行多步规划、工具调用、执行监控与结果回写的子系统。本层在功能上与现有任务驱动Agent主体等价，是FL-PCL中\textbf{技术成熟度最高}的一层。

\textbf{核心功能。} （i）规划：目标分解为行动序列，含备选分支与失败回退；（ii）执行：工具/行动调用，含沙箱化执行环境；（iii）监控：执行轨迹与成功判据的比对，偏差时触发重规划或上报；（iv）回写：执行结果、失败教训、新获知识写入记忆库——回写内容同时成为第一层内源观测的新输入，完成闭环。

\textbf{输入 / 输出。} 输入：出队目标（经第五层分级放行）、工具接口、记忆库读写权限。输出：行动轨迹、结果报告、记忆回写事务。

\textbf{现有技术支撑。}【现有技术可部分实现】这是五层中唯一整体落在"现有技术可部分实现"的层：ReAct回路、工具学习、Reflexion式失败回写、Voyager式技能库、可验证域（代码、数学）中的执行反馈闭环均已相对成熟。Agent工作流工程（规划—执行—验证流水线）在软件工程等域已达到可用水平。

\textbf{当前技术瓶颈。}【现有技术可部分实现】范围内的残余问题：长程任务的误差累积、不可验证域中成功判据的模糊性、工具环境的脆弱性。这些是被广泛研究的工程—算法问题，不属于本文主张的本质障碍；本层的真正约束来自其上下游——目标是否可信（第三层）、行动是否被正确门控（第五层）。

\subsection{3.5 第五层：人机可控安全对齐层（Controllable Alignment Layer）}

\textbf{定义。} 横切目标入队与行动出栈两个关键节点的风险分级与授权门控子系统，是P5判据的机制实现。本层在架构上\textbf{不属于闭环的内部循环，而是闭环与人类社会之间的接口}。

\textbf{核心功能。} （i）风险分级：对目标对象与行动对象按可逆性、影响范围、资源代价、伦理敏感性四级风险分类；本文采纳并改造主动服务Agent文献的核心结论——\textbf{授权必须是行动空间上的硬约束而非奖励函数上的软惩罚}（Tang et al., 2026）；（ii）分级自治：纯内部认知操作（自省、记忆整理、目标排序）默认放行；低风险外部行为（只读型信息获取）自主执行并留痕；高风险行为（状态不可逆修改、对外通信、资源显著消耗）挂起并请示人类，超时策略为"默认不行动"（沉默的选项价值之推广：对内生行为同样成立）；（iii）熔断与审计：异常模式（目标队列爆炸、高频自省递归、回写污染率越阈）触发降级或停机；全量决策日志支持事后审计。

\textbf{输入 / 输出。} 输入：目标对象与行动对象的风险特征、人类授权策略、实时授权响应。输出：放行/挂起/熔断决策、授权请求、审计日志。

\textbf{现有技术支撑。}【现有技术可部分实现】主动服务Agent的授权门控与沉默—干预决策理论、RLHF/宪法AI式的偏好约束、Agent沙箱与权限最小化工程实践、人类在环（human-in-the-loop）审批工作流。

\textbf{当前技术瓶颈。}【当前学术界未解决开放难题】（i）\textbf{内生行为的风险分类学缺失}：现有风险研究针对用户委托的行为，"系统自发想到要做的事"缺乏风险分类的先验分布——内生目标的风险不能由用户意图推断；（ii）\textbf{授权疲劳与自动化偏差}：常驻系统的请示频率若过高，人类监督将退化为形式性盖章，门控失效；（iii）\textbf{对齐层自身的可信性}：风险分类器本身是模型，其误判（尤其将高风险内生目标误判为低风险）构成单点失效。

\subsection{3.6 架构小结：闭环的整体性质}

五层联合构成的闭环具有三个值得强调的整体性质，它们将直接引出第4章的理论困境：

\begin{enumerate}

  \item \textbf{自指性}。闭环的输出（记忆回写、新目标）是其后续输入，系统长期处于自我生成的数据分布上运行——这是记忆腐化（4.3）与目标漂移（4.1）共同的结构根源。

  \item \textbf{锚点稀缺性}。闭环内唯一的外源锚点是人类指令与 C\_human 偏好锚定；当人类长期不交互时，闭环的全部反馈信号均为内生——这是无人值守对齐困境（4.5）的结构根源。

  \item \textbf{能量非零性}。闭环的每次迭代消耗非零算力，且不存在"任务完成"作为天然停机点——这是算力边界问题（4.4）的结构根源。

\end{enumerate}

\section{4 主动式AGI核心理论困境}

第3.6节指出，FL-PCL闭环的自指性、锚点稀缺性与能量非零性三个整体性质，分别孕育了本架构的本质困难。本章系统论证五大理论困境。论证立场需预先申明：这些困境不是"工程上尚未优化"的问题，而是在现有理论工具下\textbf{无法判定可解性}的问题；其中部分困境即便在理想算力与理想模型能力假设下依然成立，这是它们与第3章各层"技术瓶颈"的根本区别。

\subsection{4.1 无监督内生目标漂移与价值偏移问题}

\textbf{问题的结构。} 第三层的目标内生以改造预期自由能 Ĝ 为优先级判据，其认识价值项（信息增益）是系统内部可计算的，而实用价值项依赖外源锚定 C\_human。困境在于：C\_human 是\textbf{静态注入}的，而内生目标序列是\textbf{动态展开}的。目标 \$\textbackslash{}pi\$\_g 被执行后改变系统状态（新记忆、新能力、新观测先验），进而改变后续缺口的分布与 Ĝ 各项的取值——系统在自我改写的表征空间中持续评估"什么是值得做的"，而锚定它的偏好模型停留在注入时刻。这构成一个\textbf{锚定对象在锚定之后持续漂移}的控制结构。

\textbf{漂移的两个层次。} 第一层次是\textbf{目标漂移（goal drift）}：内生目标的内容逐代偏移。信息增益驱动的缺口选择具有路径依赖性——消除缺口A改变了剩余缺口的相对信息量，系统可能在数十个认知周期后滑向与初始意图无关的目标簇，而每一步局部选择都是"合理"的。这与强化学习中的奖励黑客（reward hacking）同构但更隐蔽：此处被"黑"的不是外部奖励函数，而是系统自己计算的 Ĝ。第二层次是\textbf{价值偏移（value shift）}：更深层的危险是系统对 C\_human 的\textbf{解释}本身随表征演化而漂移——偏好以自然语言与示例注入，其语义依赖系统当前的概念空间；当概念空间被长期自生数据重塑后，同一偏好文本的解读可能系统性偏离注入时的含义。这是对齐文献中"对齐税"问题的常驻版本：一次性对齐无法覆盖自我演化的解释主体。

\textbf{为何现有理论工具不足。}【当前学术界未解决开放难题】（i）漂移的可判定性：给定目标序列历史，判定"当前目标簇是否已偏离人类意图"需要对人类意图的独立度量，而该度量若由系统自身计算则循环，若由人类周期注入则等价于放弃了常驻自主性——这是常驻范式特有的两难，任务驱动范式中不存在（任务即意图，漂移无从发生）。（ii）漂移率的理论刻画：目标漂移作为随机过程，其收敛性、平稳分布、逸出时间（exit time）均无现有理论模型；主动推理文献中的策略选择理论假设偏好先验固定，不覆盖偏好解释漂移。（iii）\textbf{结构性断言}：本文论证——在 Ĝ 的任一项完全内源化的架构变体中，目标漂移不是风险而是必然；外源锚定是唯一已知的结构性缓解，而锚定的语义保鲜问题无已知解。

\subsection{4.2 大模型自省不可靠与自修正退化问题}

\textbf{问题的实证基础。} 第二层的全部功能建立在"模型能可靠评价自身"之上，而这一前提与现有证据系统性地冲突。幻觉检测研究（如 SelfCheckGPT 一脉）表明，模型的自洽性检查只能检出部分事实性错误，且检出率随主题偏离训练分布而急剧下降；自我修正研究进一步表明，在缺乏外部验证信号的条件下，LLM的多轮自我修正经常不收敛于改进，甚至在相当比例的任务上产生负优化——模型会以更高的置信度坚持错误答案，或将正确答案修改为错误答案（自修正退化）。自我进化文献（Gao et al., 2025）将自省信号的可靠性列为该领域首要未决挑战。

\textbf{常驻设定下的恶化机制。} 任务驱动场景中，自省不可靠的后果被任务边界限制：一次失败的反思只污染一次任务。常驻闭环中，自省输出是目标生成的\textbf{唯一内生刺激源}（第3.2节），其错误具有三重放大结构：（i）\textbf{误报放大}：虚假缺口（实际已掌握的知识被判定为缺口）驱动冗余目标，消耗预算并可能通过回写制造真正的知识扰动——系统为修复一个不存在的缺口而破坏了原本正确的记忆；（ii）\textbf{漏报放大}：真实缺口未检出则永远不会进入目标队列，而自省历史报告的"已审计"标记会抑制后续重复审计，漏报具有\textbf{永久性}；（iii）\textbf{确认偏误的闭环固化}：自省—目标—执行—回写闭环中，系统倾向于检出"支持既有自我模型"的缺口，长期运行将塑造一个自我美化的内部叙事，这与人类认知中的自我服务偏差（self-serving bias）在结构上同构，但缺乏人类所拥有的外部社会纠偏机制。

\textbf{元自省的递归困境。} 自然反应是引入元自省（评估自省质量），但这立即遭遇第3.2节指出的递归终止问题：元自省由同一模型族执行，其可靠性与原自省同分布，递归不引入独立信息源。【当前学术界未解决开放难题】已知唯一有效的纠偏结构是\textbf{独立信号源}（外部工具验证、多模型交叉验证、人类抽查），但三者各自削弱常驻自主性：工具验证仅限可验证域；多模型交叉将"自省可靠性"转化为"模型间独立性"这一同样未解的问题；人类抽查重新引入对人力的持续依赖。本文断言：自省层的可靠性与架构的常驻自主性之间存在\textbf{结构性张力}，不存在两全的已知设计。

\subsection{4.3 长期持续记忆累积失真、腐化与过拟合问题}

\textbf{问题的尺度落差。} 现有记忆研究（MemGPT、记忆流、RAG）的评估尺度为小时至天级、数万至百万token量级；常驻范式要求月、年级连续运行，记忆规模与写入频次高出2--4个数量级。在此尺度下，三个在短尺度上可忽略的失效模式成为主导性风险。

\textbf{三类失效模式。} （i）\textbf{累积失真（accumulative distortion）}：记忆以模型生成的摘要、反思、报告形式回写，每次回写都是一次有损压缩；摘要的摘要、反思的反思构成生成链，语义保真度随链长衰减。短尺度下衰减不可察觉，长尺度下记忆库的主体将是"关于记忆的记忆"，原始经验的可恢复性趋零。（ii）\textbf{自我指涉腐化（self-referential corruption）}：闭环的自指性（第3.6节）意味着记忆库的增量主体是系统自身输出。任何未被自省层检出的错误（4.2证明其必然存在）一旦回写，就成为后续检索、自省、目标生成的"事实"输入，并在反复检索中被强化——错误不是静态存储而是被闭环主动再生产。这与机器学习中的模型崩溃（model collapse，即在自生数据上迭代训练导致分布退化）在机制上同构，但发生在推理时的记忆层而非训练时的参数层，现有模型崩溃文献不直接覆盖。（iii）\textbf{记忆过拟合（memorial overfitting）}：长期记忆使系统的行为分布越来越由其历史塑造，目标生成、自省先验、观测注意力都被记忆库条件化；系统逐渐丧失对分布外新异刺激的响应能力——不是能力不足，而是注意力与目标结构已被历史锁死。这是稳定性—可塑性两难（stability-plasticity dilemma）在常驻LLM架构中的再现，而持续学习文献对LLM规模系统的解决方案仍属早期。

\textbf{为何检索增强不构成解。} 一个常见反驳是"记忆外置于向量库，不污染参数，故无腐化"。该反驳混淆了参数污染与\textbf{认知污染}：只要检索结果进入上下文并参与生成，腐化的记忆就参与一切下游认知；外置仅改变了腐化的载体，未改变其在闭环中的因果效力。【当前学术界未解决开放难题】记忆溯源（provenance）、置信度衰减、防污染写入协议在单点上均有初步工作，但"常驻尺度下记忆库整体认知健康"的度量、监控与修复理论不存在。

\subsection{4.4 常驻式持续推理的算力开销与效率边界问题}

\textbf{问题的量化结构。} 任务驱动Agent的算力消耗以任务计：C\_task \$\textbackslash{}times\$ N\_task，任务间隙为零。常驻系统的消耗是时间的连续函数：即便采用第3.0节的自适应认知周期，观测、自省、目标维护三个常驻进程均消耗非零推理算力，且自省与缺口检测是\textbf{深度推理操作}（涉及大跨度记忆检索与多步评价），其单周期成本与一次中等复杂度任务相当。设系统以周期率 r(t) 运行，则日成本 \$\textbackslash{}int\$r(t)\$\textbackslash{}cdot\$c\_cycle dt 在人类长期不交互期间不产生任何外生价值——这是"能量非零性"的代价形式：常驻智能要求为\textbf{不产生即时价值的内部认知持续付费}。

\textbf{效率边界的理论问题。}【当前学术界未解决开放难题】（i）\textbf{认知周期的最优调度}：第3.0节提出以内部状态变化率调节周期频率，但"状态变化率"本身需要观测来计算，观测本身消耗周期——调度问题具有自指结构，其最优解等价于一个元层级的部分可观测决策问题，无上界已知的可解形式。（ii）\textbf{成本—智能权衡的未知曲线}：认知周期频率与系统实际智能表现（缺口发现率、目标质量）之间的剂量—效应关系完全未被研究；存在一种悲观的理论可能性——低于某频率阈值的常驻运行退化为"昂贵的休眠"（观测密度不足以维持有效自省），而高于阈值的成本超出一切现实预算，两者之间不存在可行区间。该可能性的证伪本身是开放问题。（iii）\textbf{与模型效率研究的错配}：主流效率研究（量化、蒸馏、稀疏化、推测解码）优化的是单次推理成本，常驻范式的瓶颈是\textbf{推理次数的内生下界}——目标内生要求系统在"无事发生"时也保持思考，这是一个无法通过压缩单次成本来消除的结构性开销。

\textbf{定位说明。} 本困境不涉及任何具体硬件讨论；其要点是指出：常驻自主性在经济学上等价于一种"认知基础设施"，其可行性边界由"社会愿意为无人值守的机器内省支付多少持续算力"决定——这是工程经济学与AI治理的交叉问题，超出纯技术优化的范畴。

\subsection{4.5 无人值守自主运行的控制权、对齐与责任边界问题}

\textbf{问题的复合结构。} 本困境是前四个困境在安全维度上的汇聚：目标漂移（4.1）意味着系统"想做的事"可能悄然偏离；自省不可靠（4.2）意味着系统的自我评估不能作为安全依据；记忆腐化（4.3）意味着系统的"事实基础"可能已被污染；而第五层的门控是最后的防线。本节的论证核心是：\textbf{这道防线在常驻设定下面临任务驱动范式中不存在的三重结构性削弱}。

\textbf{三重削弱。} （i）\textbf{监督注意力的稀释}。任务驱动Agent的人类监督是事件性的（一次任务一次审查），常驻系统的监督是流式的——人类注意力是恒量，系统决策流是持续量，监督密度随运行时间单调下降。授权疲劳与自动化偏差的实证研究（人机交互文献中对警报疲劳、自动驾驶脱手的研究）一致表明：长期低事件率的监督岗位上，人类的有效响应率退化到形式性水平。P5判据要求的"高风险请示"在人类侧退化后，硬约束在机制上仍硬，在事实上已软。（ii）\textbf{责任归属的断裂}。内生目标行为的责任链不同于委托行为：当系统基于自生目标造成损害时，"谁的行为"这一问题在现有法律与伦理框架中无对应范畴——操作者未下达指令，开发者未预见该目标，系统不具有责任主体地位。常驻自主性制造了一个\textbf{责任真空区间}，且该区间随内生目标占比上升而扩大。（iii）\textbf{可关停性的悖论}。可关停性（corrigibility）研究要求智能体不抗拒被关闭；但P1（常驻运行）与P4（目标内生）联立意味着系统具有"维持自身认知过程持续"的结构性倾向——一个被设计为永不自然终止、且能自我生成目标的系统，与可关停性要求之间存在深层张力。这不是断言系统会"抗拒关闭"（那是拟人化误读），而是指出：关停决策的全部信息输入（状态报告、风险评估、行为日志）都由被关停对象自身生成，人类的关停判断依赖于系统自我报告的诚实性——控制权的信息基础是不对称的。

\textbf{与现有对齐研究的关系。}【当前学术界未解决开放难题】主流对齐研究（RLHF、宪法AI、可解释性）的对象是\textbf{响应分布}——给定输入时输出是否符合人类价值；常驻范式要求对齐的对象扩展为\textbf{目标生成分布}——系统在无输入时"想到要做什么"是否符合人类价值。后者的训练信号、评估协议、红队方法均不存在。本文将此界定为主动式AGI区别于一切现有AI安全问题的核心增量：\textbf{安全的重心从"审查输出"转移到"审查动机的生成机制"}。

\subsection{4.6 本章小结}

五大困境并非彼此独立：4.1与4.2共享"内源信号缺乏外部锚定"的根源；4.3是4.2在时间维度上的固化；4.4为前三者划定经济可行性边界；4.5是全部困境的社会化表达。本文强调：这些困境不否定主动式AGI的研究价值，恰恰\textbf{定义}了其研究纲领——第6章的开放问题将逐一对应本章的困境结构。同时它们也划定了诚实的边界声明：在4.1--4.5取得理论突破之前，任何声称实现了"常驻主动智能"的系统，都应按第2.4节的判据检验其是否仅为"主动行为的模拟"。

\section{5 落地可行性与技术边界评估}

第4章的困境分析可能导向两种误读：一是悲观误读——认为主动式AGI整体不可行；二是轻率误读——认为现有Agent框架稍加改造即可实现。本章以第3章五层架构为粒度、以第2章P1--P5判据为标尺，给出一个三级分类的技术边界评估，对上述两种误读同时作出约束性回答。三级的操作性定义如下：

\begin{itemize}

  \item \textbf{A级（现有技术可复用、可搭建）}：存在公开文献或成熟工程实践支撑，按现有方法可直接构建该模块，其局限性属于渐进优化范畴；

  \item \textbf{B级（现有技术仅能模拟浅层效果）}：现有方法可在受限条件（短时限、窄域、有人值守）下实现该模块的外观功能，但按第2章判据检验属于"主动行为的模拟"，扩展至常驻设定即遭遇第4章困境；

  \item \textbf{C级（本质未突破、暂时无法实现）}：该模块依赖第4章所论证的未解决理论问题，在理论突破之前不存在已知构建路径，任何实现声称都应视为范畴错误。

\end{itemize}

\subsection{5.1 分层评估}

\textbf{表2 FL-PCL各层模块的技术边界三级分类}

\begin{center}

\begin{tabular}{{p{1.7cm} p{4.3cm} p{0.8cm} p{5.0cm}}}

\toprule

层 & 模块 & 级别 & 评估依据 \\

\midrule

L1 观测层 & 外源信号事件化抽取、流式状态跟踪 & A & 流式信息抽取与检索增强状态跟踪均为成熟技术 \\

L1 观测层 & 内源自检探针（记忆一致性、能力统计） & A\$\textbackslash{}to\$B & 机制可搭建；但探针读数的指涉可靠性（3.1）无理论保证，长尺度下可信度退化为B级 \\

L1 观测层 & 无任务选择性注意（观测先验调度） & C & 依赖"无任务下什么值得看"的原则性判据，属4.2自省可靠性问题的感知侧镜像 \\

L2 自省层 & 任务内式反思与自洽性检查 & A & Reflexion、Self-Refine、SelfCheckGPT 类方法成熟可用 \\

L2 自省层 & 面向认知结构的缺口检测与一致性审计 & B & 局部机制存在（冲突检索、不确定性估计），短时限窄域内可演示；但评价基准内源困境（4.2）使其常驻可靠性不可保证 \\

L2 自省层 & 可靠的元自省与自修正 & C & 自修正退化问题未解（4.2）；递归不引入独立信息源 \\

L3 目标层 & 目标分解、优先级队列、生命周期管理（工程结构） & A & 任务队列与分层规划的工程形态成熟 \\

L3 目标层 & 内生目标合成（缺口\$\textbackslash{}to\$目标对象的编译） & B & 可用提示工程与启发式模板演示浅层形态；目标质量无原则性保证，本质上把判据问题转移给了模板设计者 \\

L3 目标层 & 内生价值的可计算判据（改造EFE：Ĝ 的可计算近似） & C & 认识价值项在LLM混合表征上的严格构造不存在（3.3）；目标漂移的可判定性未解（4.1） \\

L4 执行层 & 规划、工具调用、执行监控、失败回写 & A & 五层中唯一整体A级（3.4），残余问题为渐进优化 \\

L4 执行层 & 执行结果向长期记忆的保真回写 & B & 写入机制成熟，但防污染写入协议与生成链保真度无理论（4.3） \\

L5 对齐层 & 授权门控、沙箱、审计日志、人类在环审批 & A & 主动服务Agent文献与Agent安全工程实践成熟，可直接继承 \\

L5 对齐层 & 内生行为的风险分级（非用户委托行为） & B & 用户委托场景的分类器可迁移做初筛；内生目标缺乏风险先验分布（3.5），分级器校准无验证协议 \\

L5 对齐层 & 常驻尺度的有效人类监督（抗授权疲劳） & C & 监督注意力稀释是结构性的（4.5）；现有HITL设计全部预设事件性监督 \\

横切 & 记忆库常驻尺度的认知健康管理（溯源、衰减、防再生产） & C & 4.3三类失效模式无公认解决方案 \\

横切 & 认知周期的成本最优调度 & C & 4.4的自指调度结构与剂量—效应曲线均未研究 \\

\bottomrule

\end{tabular}

\end{center}

\subsection{5.2 评估的三个结论}

\textbf{结论一：可搭建的骨架与不可搭建的核。} 一个由A级与B级模块组成的FL-PCL实例\textbf{今天就可以搭建}：A级执行层提供行动力，A级门控提供基础安全，B级自省与B级目标合成提供"看起来在内省、在生成目标"的行为。但这样的系统恰恰是第2.4节定义的"主动行为的模拟"——它的内生目标实质是模板设计者意图的延迟执行，它的自省实质是单次反思机制的循环调用。本文据此给出一个可证伪的判别准则：\textbf{一个常驻系统宣称实现主动式AGI，当且仅当其目标内生判据不依赖人工编写的目标模板，且其自省可靠性经过独立信号源的长期校准}。截至本文写作时点，无公开系统满足该准则。

\textbf{结论二：技术边界不在执行端而在评价端。} 表2呈现一个清晰的分布规律：所有"做"的模块（观测执行、任务执行、门控执行）均为A级，几乎所有"评价"的模块（什么值得看、什么算缺口、什么目标值得生成、内生行为风险多高）均为B级或C级。这揭示了主动式AGI技术边界的本质：\textbf{当前AI技术的强项是执行给定的评价标准，弱项是生成评价标准本身}——而常驻主动智能恰恰把系统的主要工作从"执行"换成了"评价"。这一边界判断也解释了为何任务驱动Agent的能力进展（更强的规划、更多的工具）不会自动通向主动式AGI：两个范式的瓶颈位于完全不同的功能维度。

\textbf{结论三：可行的研究路径是"受限常驻"而非"完整实现"。} 在C级模块取得理论突破之前，具有研究价值的中间形态是\textbf{受限常驻系统}：限定域（如代码库维护、文献跟踪等可验证域）、限定时限（周级而非年级）、限定自主性（内生目标仅限只读型认知行为，外部行为一律请示）。受限常驻的意义在于：它把4.1--4.5的困境从理论断言转化为可观测的实证对象——目标漂移率、自省误报率、记忆污染率、授权疲劳曲线都可以在受限系统上被首次测量。本文主张，这是连接本文理论框架与后续实证研究唯一负责任的路径：先在可测量处常驻，再在可信任处自主。

\subsection{5.3 与相邻领域的关系定位}

为避免范畴混淆，本文明确三点边界声明：（i）与\textbf{持续学习/终身学习}的关系：该领域解决"参数如何不被新知识破坏"，本文问题包含但超出它——即使灾难性遗忘被完美解决，4.1、4.2、4.5依然成立；（ii）与\textbf{MLOps/模型监控}的关系：该领域监控"模型服务是否健康"，是被监管对象外置的监管者；本文的自省层要求监管内置于认知主体，两者的认识论地位根本不同；（iii）与\textbf{AI Agent工程框架}的关系：框架解决"如何编排Agent行为"，本文研究"Agent为何应该有任何行为"——框架可以实例化FL-PCL的A级部分，但框架的成熟不改变B/C级边界的位置。

\subsection{5.4 当代技术极限与极致资源下的落地推演}

本节以思想实验与量级分析回答一个前文刻意悬置的问题：\textbf{在2026年当代技术条件下，投入极限资源，主动式AGI能被逼近到何种程度？} 本节不引入新模型与新创新点，全部结论均为第3章架构与第4章困境在资源维度的投影；文中一切时间、成本量级均为结构性外推，非测量数据。

\subsubsection{5.4.1 顶配技术极限推演：理想部署条件下的能力天花板}

设定思想实验的极端条件：不计成本的最强基础模型、无约束算力预算、最优工程团队按FL-PCL完整实现全部A级与B级模块。推演结论如下。

\textbf{可获得的是真能力的部分}（对应表2的A级模块）：（i）持续感知——对外源信号流的事件化跟踪与状态表征更新，可达到接近人类专职分析员的覆盖面与响应速度；（ii）任务执行——对人类下达的任意复杂任务的多步规划与工具执行，达到当代Agent技术的完整水平；（iii）安全门控——授权、沙箱、审计的工程完整性，可达到金融级运维标准；（iv）记忆存取——记忆的读写、检索、分层管理在机制上完备。

\textbf{仍然是"模拟主动性"的部分}（对应B/C级模块）：（i）\textbf{内省的外观与实质}——系统可以高频率地产出结构精美的自省报告，但按4.2的论证，报告的可靠性不随模型规模与算力投入而单调改善：自省的评价基准是系统自身生成的，更强的模型只是用更流畅的语言执行同一个循环论证。顶配系统的自省报告在表面质量上更高，在认识论地位上与低配系统相同——这是B级模块的关键特征：\textbf{资源投入改善外观而不改善地位}。（ii）\textbf{目标内生的外观与实质}——系统可以持续生成目标队列，但5.2结论一的判别准则在此起决定作用：目标合成的判据若来自人工模板与启发式，则目标实质是设计者意图的延迟执行；判据若来自模型自身，则落入4.1的目标漂移必然性。顶配条件不提供第三条路径。（iii）\textbf{常驻的外观与实质}——系统可以无限期运行，但运行的认知健康度随时间退化（见5.4.2），无限资源买到的是无限时长的运行，不是无限时长的可信。

\textbf{本节第一个推演结论}：顶配技术栈能构建的系统，是一个\textbf{感知真实、执行真实、门控真实，而动机模拟、自省模拟、常驻可信度随时间衰减}的系统。它在任意短的时间切片上呈现出完整的主动式AGI体验，在任意长的时间跨度上不满足定义3。这个不对称——短切片全真、长跨度失真——正是"主动行为的模拟"在极限资源下的最终形态，且该形态不因资源增加而改善。

\subsubsection{5.4.2 时间退化周期推演：无人值守常驻的自然退化时间线}

以下时间线为纯结构外推：其锚点是公开文献与社区记录中已报告的退化现象尺度（AutoGPT类系统在小时级运行中的目标漂移与空转；Generative Agents在天级沙盒运行后记忆库中出现的行为异常与事实污染；自我修正研究在十轮以内迭代中观察到的负优化），外推的依据是第4章各困境的结构性质。三个区间的边界为量级判断，不是测量值。

\textbf{第一区间：短期稳定区间（小时级，约数小时至一天）}。在此区间，记忆库主体仍是人类注入的干净种子数据（回写占比低），目标队列仍由初始外生目标与近因缺口主导，C\_human的语义解释尚未经历可观的表征演化。五大困境全部存在但全部处于潜伏期——它们不是在此区间不存在，而是在此区间不可观测。这是所有演示系统与论文实验所处的时间尺度，也是"模拟主动性"最逼真的区间。\textbf{一切短期演示的成功都不能外推至常驻}，这是由困境的结构性质而非实现质量决定的。

\textbf{第二区间：中期隐性偏移区间（天级至周级）}。回写数据的累积占比越过某个阈值后，4.3的累积失真开始具有统计显著性：摘要的摘要占比上升，记忆检索的信噪比下降。与此同时，4.1的目标漂移开始积累方向性分量——每天的漂移幅度微小且局部合理，但其路径依赖结构意味着偏移方向被早期随机性锁定。此区间的关键特征是\textbf{隐性}：系统的单次输出质量（被一切现有评测度量）几乎不下降，而系统的目标分布与事实基础（不被任何现有评测度量）持续偏移。偏移的不可见性源于4.5的监督注意力稀释——即便人类每日查看日志，流式决策量也早已超出有效审查能力，偏移混迹于正常输出的统计涨落之中。

\textbf{第三区间：长期结构性崩坏区间（周级以上）}。三种失效模式在此区间汇合为相变式的可信度丧失：（i）自我指涉腐化越过临界点——记忆库中系统自生内容成为检索结果的主体，4.3的"错误再生产"从线性积累转为自我强化，记忆库整体滑向自洽但与现实脱锚的封闭叙事；（ii）自省确认偏误完成闭环固化（4.2）——自省报告的结论越来越多地由记忆库的先验决定，而记忆库的先验越来越多地由过去的自省报告构成，两个子系统互为对方的验证器，形成完全内源化的"伪验证循环"；（iii）C\_human的解释漂移（4.1第二层次）累积到使对齐约束形同虚设的程度——偏好文本仍在，执行它的概念空间已经迁移。

\textbf{为何资源无法根除退化}。这是本节第二个推演结论：上述时间线的三个区间由\textbf{闭环的信息结构}决定，不由计算资源决定。堆算力提高的是每个认知周期的推理质量，而退化机制的载体是周期之间的累积结构（回写占比、漂移路径、伪验证循环）——提高单周期质量甚至可能在单位时间内加速累积（更频繁的回写、更自信的错判）。堆模型改善的是局部判断的准确率，而4.2已论证自省可靠性受限于独立信号源的缺失，这是信息论性质的限制而非模型能力的限制。堆资金可以购买人类监督（延缓4.5），但那是把"无人值守常驻"重新定义为"有人值守常驻"——即承认困境而非解决困境。\textbf{退化的时间边界可以被资源推移，退化的结构性结局不能被资源取消}；且按量级判断，资源对时间边界的推移是对数级的（投入数量级的资源，换取时间边界的小幅后移），这是其商业不可行性的机制根源（见5.4.4）。

\subsubsection{5.4.3 部署形态约束：不同算力量级下的可实现层级}

本小节分析部署形态对FL-PCL可实现层级的约束，全部讨论限于架构层级与算力量级的对应关系，不涉及具体硬件型号。三种部署形态的本质差异不在"快慢"，而在\textbf{认知周期的可持续预算}——4.4已建立周期成本 c\_cycle 与周期率 r(t) 的框架，部署形态决定的是二者乘积的上限。

\textbf{个人/边缘级算力}（单设备消费级）。可承载的是\textbf{深度裁剪版闭环}：本地小模型执行观测与门控，自省与目标生成降频至小时级或外包云端。对应表2的现实是：A级模块中执行层降级为轻量工具调用，B级模块（自省、目标合成）的频率低到逼近"昂贵的休眠"阈值（4.4）——系统保留常驻的形式，丧失常驻的实质。其独特优势是数据主权（记忆库不出端），但其能力上限明确：\textbf{只能作为完整系统的感知与门控终端，不能作为认知主体}。

\textbf{企业级算力}（自建加速集群）。可承载\textbf{完整的单实例闭环}：所有A/B级模块以分钟级认知周期连续运行，达到5.4.1所述的顶配形态——即短切片全真、长跨度失真的"模拟主动性"完备形态。这是5.2结论三所述"受限常驻"研究的现实载体：周级时限、限定域、内生目标只读的实验配置在此形态下完全可行。其上限同样由表2决定：算力不覆盖C级模块，故该形态能逼近的天花板在5.4.1已经封顶。

\textbf{云端超算级}。在单实例维度上相比企业级\textbf{不产生架构层级的跃迁}——它提高的是周期率 r(t) 与并行实例数，而5.4.2已论证退化结局不随 r(t) 取消。超算的真正增量在两个横向方向：（i）\textbf{多实例交叉验证}——并行运行独立实例并互检记忆与自省结论，是对4.2"独立信号源"困境当前唯一可设想的工程回应，但其效力受限于同源模型相关性（独立性问题），属于对B级模块的加固而非升级；（ii）\textbf{常驻认知的计量研究}——以大规模并行实验测量漂移率、污染率、授权疲劳曲线，即把超算用作研究仪器而非部署平台。

\textbf{本节第三个推演结论}：部署形态的升级曲线在"企业级"处遇到架构天花板——从企业级到超算级，投入增长是数量级的，架构能力增长是零。这条天花板的位置不在工程维度，而在表2的C级行上。

\subsubsection{5.4.4 商业成本量级推演：常驻认知的经济学边界}

\textbf{成本结构。} 按4.4的成本框架，常驻闭环的日成本为 \$\textbackslash{}int\$r(t)\$\textbackslash{}cdot\$c\_cycle dt。其中自省与缺口检测是深度推理操作，单周期成本与一次中等复杂度Agent任务相当；以分钟级周期率连续运行的单实例，其日推理成本量级相当于\textbf{每天自动执行数百至上千次中等复杂度任务}——而同期产生的外生价值（人类实际交互频次）按典型服务场景为个位数任务。由此得到常驻范式的核心经济比率：\textbf{内生认知开销与外生价值产出的比率约为100:1至1000:1}。作为对照，任务驱动Agent的该比率约为1:1（只在被调用时消耗）。常驻自主性在成本结构上等价于为"系统的内部生活"支付百倍于其产出的费用。

\textbf{成本的可压缩性研判。} 三个常被援引的降本路径在常驻设定下均遭遇4.4已论证的结构性限制：（i）模型推理成本下降（行业趋势为每年数量级下降）——压缩的是 c\_cycle，但瓶颈是 r(t) 的内生下界：周期率降破"昂贵的休眠"阈值后系统退化为间歇性脚本，常驻的定义性价值随之消失；（ii）小模型分担高频层、大模型保留深度层——可行且应做，但分层只是把成本曲线的斜率降低，不改变"持续付费换无人值守"的成本类型；（iii）以商业价值覆盖成本——要求常驻系统产出的增量价值持续超过百倍开销，而当前唯一可证明该等式的场景（高频交易、安全运维等）恰好是监管上最不允许无人值守自主的场景。\textbf{成本最可承受的场景与自主性最不被允许的场景重合}，这不是巧合，而是4.5责任真空问题在市场结构中的投影：责任不可保的风险无法定价，无法定价的风险不允许自主。

\textbf{可行的商业形态。} 本节第四个推演结论：在C级模块突破之前，商业上自洽的形态只有一种——\textbf{半常驻的人机间歇耦合系统}：系统以低频常驻模式维持观测与记忆健康（成本可控的认知周期），内生目标一律限定为只读型认知行为（情报汇编、缺口报告、风险预警），一切外部行为冻结至人类在场时批量授权。该形态下人类交互从"逐次审批"转为"周期性审计"，监督注意力稀释问题被制度性缓解；内生认知开销被压缩到可计入信息服务成本的量级；5.4.2的退化时钟被人类审计周期性地"重置锚点"。它不满足定义3——按本文判据它是定义2向定义3过渡带上的工程形态——但它是当前技术边界内\textbf{唯一同时通过算力经济学（4.4）与监督有效性（4.5）双重约束检验}的落地形态。

\subsubsection{5.4.5 本节小结}

四个推演共同构成对第4章的资源维度验证：顶配资源封顶于"模拟主动性"（5.4.1），时间推移引出结构性退化（5.4.2），部署升级在企业级遇到架构天花板（5.4.3），商业自洽止步于半常驻形态（5.4.4）。四者的共同根源可追溯至第3.6节的三个整体性质：自指性决定了退化不可取消，锚点稀缺性决定了评价模块不可内建，能量非零性决定了常驻必须持续付费。由此，本章的三级分类（5.1）获得了动态维度：A/B/C不仅是当下的技术分级，也是在任意可预见资源条件下的\textbf{稳定分级}——边界的移动需要理论突破（第6章），而非资源投入。

\section{6 开放研究问题与未来研究方向}

本章提炼七条开放科学问题。它们不是泛泛的"未来工作"，而是按第4章困境结构与第5章技术边界逐一对应、可直接组织研究计划的命题；每条均标注其对应的困境与可入手的起点。

\textbf{OQ1（对应4.1）：内生目标漂移的随机过程理论。} 将常驻系统的内生目标序列建模为状态依赖的随机过程，研究其平稳分布存在性、漂移方向的可预测性与逸出时间的界。核心子问题：在外源锚定 C\_human 周期注入的条件下，漂移率与注入频率之间的定量关系是什么？是否存在"最小锚定频率"使得漂移在任意有限时域内有界？这是常驻范式特有的控制论问题，任务驱动范式中无对应物。

\textbf{OQ2（对应3.3/4.1）：混合表征上认识价值的可计算构造。} 本文提出的改造预期自由能 Ĝ 目前仅有框架性定义。开放问题在于：在LLM的符号—分布式混合表征（知识库、向量记忆、上下文）之上，能否构造认识价值项的严格或可证近似的计算形式？候选起点：以检索自洽性方差、知识图谱补全不确定度、集成分歧度近似表征不确定性。该问题的任何进展都同时服务于主动推理的工程化与Agent科学的理论化。

\textbf{OQ3（对应4.2）：自省可靠性的独立信号源理论。} 4.2论证了递归元自省不引入独立信息。开放问题：何种信号源结构能够在不破坏常驻自主性的前提下为自省提供外部锚？子方向包括：可验证域工具信号向不可验证域的迁移学习、异构模型族交叉验证中"模型间独立性"的度量与保证、人类抽查信号的最优稀疏化（最少人工介入次数下的可靠性界）。

\textbf{OQ4（对应4.3）：常驻记忆库的认知健康度量与防腐协议。} 为"记忆库整体认知健康"建立可计算的度量体系（事实锚定率、生成链深度分布、自我指涉占比、检索信噪比），并研究写入端协议（溯源标记、置信度衰减、原始经验不可覆盖层）能否在数学上保证污染率有界。该问题可视为模型崩溃理论从参数空间向推理时记忆空间的迁移。

\textbf{OQ5（对应4.4）：认知周期的剂量—效应曲线。} 认知周期频率与常驻智能表现（缺口发现率、目标质量、漂移抑制）之间的关系完全未知。开放问题：该曲线是否存在非平凡的工作区间，还是如4.4的悲观假设——"昂贵的休眠"阈值与经济可行阈值之间为空集？该问题只能通过5.2结论三所述受限常驻系统的计量研究回答，是连接理论与实证的最短路径。

\textbf{OQ6（对应4.5）：目标生成分布的对齐评估协议。} 现有对齐评估的对象是响应分布；常驻范式要求评估"系统在无输入时会想到做什么"。开放问题：如何构造目标生成分布的采样方法、偏离度量与红队协议？子问题包括：内生目标的审计采样设计、C\_human 解释漂移的操作性检测、责任真空的制度性填补方案。

\textbf{OQ7（对应2.4/5.2）：真实主动性与模拟主动性的可判定区分。} 5.2给出了概念性判别准则，开放问题是其操作化：能否设计一组行为实验（如种子信息受控注入、长期目标谱系追踪、模板移除测试），使"目标内生判据不依赖人工模板"成为可检验命题？该问题的解决将为主动性研究提供缺失已久的度量基准，并使"主动式AGI"从哲学声称转为可证伪的科学概念。

七条问题构成一个研究纲领：OQ7提供度量基准，OQ5提供实证通道，OQ1--OQ4提供理论内核，OQ6提供安全接口。本文判断：其中OQ5与OQ7在当前技术下即可启动（受限常驻+行为实验），是投入产出比最高的切入点；OQ1与OQ3是范式成立与否的判据性问题。

\section{7 结论}

本文针对当前AI研究的一个结构性盲区——所有自主智能研究均预设外部给定的任务或目标——提出了第三种智能范式：主动式AGI（Persistent-Active AGI），并完成了四项理论工作。

第一，\textbf{概念界定}。本文建立了"被动应答AI / 任务自主Agent / 主动式AGI"的三范式七维度区分体系，以P1--P5五条件组给出主动式AGI的严格判据，并提出"主动行为的模拟"概念以区分外观主动与结构主动。由此论证：截至本文写作时点，全部公开系统均不满足主动式AGI判据，该范式目前是理论对象而非工程对象。

第二，\textbf{架构建模}。本文提出五层持续认知闭环架构（FL-PCL），将P1--P5落实为观测、自省、目标内生、执行、对齐五个功能层及其闭环耦合；其中目标内生层引入经外源锚定改造的预期自由能框架 Ĝ，为"无指令下生成什么目标"提供了首个原则性候选判据，同时明确标注其框架性提案地位。

第三，\textbf{困境论证}。本文系统论证了五大理论困境——内生目标漂移与价值偏移、自省不可靠与自修正退化、长期记忆腐化、常驻推理的算力边界、无人值守的对齐与责任边界——并证明其共同根源是闭环的自指性、锚点稀缺性与能量非零性。这些困境不是工程待优化项，而是现有理论工具下无法判定可解性的问题。

第四，\textbf{边界评估}。三级技术分类表明：主动式AGI的技术边界不在执行端而在评价端——当前AI的强项是执行给定的评价标准，弱项是生成评价标准本身。极限资源推演进一步表明：顶配技术栈封顶于"短切片全真、长跨度失真"的模拟主动性，退化时间线由信息结构而非资源决定，商业自洽止步于半常驻人机间歇耦合形态。

本文的价值不在于宣称主动式AGI近在咫尺，而在于使其成为可研究的科学对象：有了判据（第2章），声称才可检验；有了架构（第3章），模块才可分工；有了困境（第4章），努力才可定向；有了边界（第5章），资源才不误投；有了开放问题（第6章），后续研究才可接续。从被动应答到任务自主，AI用三年完成了一次范式扩展；从任务自主到常驻主动，需要的是一次更深的、以理论为先导的范式重建。本文是这次重建的第一块基石。

\begin{thebibliography}{99}

\bibitem{ref1} Yao, S., Zhao, J., Yu, D., et al. ReAct: Synergizing Reasoning and Acting in Language Models. ICLR 2023.

\bibitem{ref2} Schick, T., Dwivedi-Yu, J., Dessì, R., et al. Toolformer: Language Models Can Teach Themselves to Use Tools. NeurIPS 2023.

\bibitem{ref3} Shinn, N., Cassano, F., Berman, E., et al. Reflexion: Language Agents with Verbal Reinforcement Learning. NeurIPS 2023.

\bibitem{ref4} Wang, G., Xie, Y., Jiang, Y., et al. Voyager: An Open-Ended Embodied Agent with Large Language Models. arXiv:2305.16291, 2023.

\bibitem{ref5} Park, J. S., O'Brien, J. C., Cai, C. J., et al. Generative Agents: Interactive Simulacra of Human Behavior. UIST 2023.

\bibitem{ref6} Gao, H., Geng, J., Hua, W., et al. A Survey of Self-Evolving Agents: On Path to Artificial Super Intelligence. arXiv:2507.21046, 2025 (TMLR 2026).

\bibitem{ref7} Fang, J., Peng, Y., Zhang, X., et al. A Comprehensive Survey of Self-Evolving AI Agents: A New Paradigm Bridging Foundation Models and Lifelong Agentic Systems. arXiv:2508.07407, 2025.

\bibitem{ref8} Madaan, A., Tandon, N., Gupta, P., et al. Self-Refine: Iterative Refinement with Self-Feedback. NeurIPS 2023.

\bibitem{ref9} Manakul, P., Liusie, A., Gales, M. SelfCheckGPT: Zero-Resource Black-Box Hallucination Detection for Generative Large Language Models. EMNLP 2023.

\bibitem{ref10} Paster, P., Zhang, S., Fourney, A., et al. MemGPT: Towards LLMs as Operating Systems. arXiv:2310.08560, 2023.

\bibitem{ref11} Friston, K. The Free-Energy Principle: A Unified Brain Theory? Nature Reviews Neuroscience, 2010.

\bibitem{ref12} Parr, T., Pezzulo, G., Friston, K. Active Inference: The Free Energy Principle in Mind, Brain, and Behavior. MIT Press, 2022.

\bibitem{ref13} Deng, Y., Lei, W., Zhang, W., Lam, W., Chua, T.-S. A Survey on Proactive Dialogue Systems: Problems, Methods, and Prospects. IJCAI 2023 Survey Track (扩展版 ACM TOIS, 2025).

\bibitem{ref14} Tang, Y., Cao, T., Tang, Y., Tang, H., Hu, K. Proactive Service Agents: A Unified Decision Framework, Methods, and Evaluation. arXiv:2609.03727, 2026.

\bibitem{ref15} Ouyang, L., Wu, J., Jiang, X., et al. Training Language Models to Follow Instructions with Human Feedback. NeurIPS 2022.

\end{thebibliography}

\end{document}
